\section{Types}
\label{sec:types}

\begin{definition}
A \textbf{type} is a specification of the operations, the data structure, the
external interface and the internal representation of the elements of a
program, its values.
\end{definition}

A type system gives a programming language a logic: rules that describe which
transformations are allowed on which data. This block follows
Watt~\cite[Ch.~2, 8]{watt2004concepts}, Friedman and
Wand~\cite[Ch.~4]{friedman2008essentials} and
Lee~\cite[Sec.~1.5]{lee2017foundations}.

\subsection{Kinds of typing}

The lecture compares three ways in which languages handle types
(Table~\ref{tab:typing-kinds}).

\begin{table}[h]
\centering
\small
\begin{tabular}{l p{0.3\textwidth} p{0.38\textwidth}}
  \toprule
  Typing & Values & When operations are checked \\
  \midrule
  Strong   & each value has one type & statically, while the program is built \\
  Inferred & each value has one type & statically, from types given by the programmer or deduced from the use of each value \\
  Dynamic  & a value can have different types at different times; at run time it is stored as a pair (value, type tag) & at run time \\
  \bottomrule
\end{tabular}
\caption{Kinds of typing.}
\label{tab:typing-kinds}
\end{table}

Discussions about type systems usually turn around four questions, and a
language answers each one independently:
\begin{itemize}
  \item static or dynamic: are types checked before the program runs or
    while it runs?
  \item strong or weak: may the language reinterpret a value as another type
    without an explicit conversion? In Python \lstinline|"3" + 3| is an
    error; in JavaScript it gives \lstinline|"33"|.
  \item explicit or inferred: does the programmer write the types, or does
    the compiler deduce them?
  \item nominal or structural: are two types compatible because of their
    names, or because of their shape?
\end{itemize}

\subsection{Basic, primitive and composite types}

Following Martin-L\"of, three types are the basis of all others. The empty
type $\bot$ has no elements; it is used to express negation, $\neg A ::= A \to
\bot$. The unit type has a single element, written \texttt{()}, \texttt{unit} or
\texttt{null} depending on the language. The two type has two values,
\texttt{true} and \texttt{false}. All other types have infinitely many elements.

Languages provide \textbf{primitive types} (integers, booleans, characters,
pointers) and let the programmer define \textbf{composite types} (functions,
lists, objects, arrays). Primitive types are defined by the range of values
they represent: \texttt{int} is $-2147483648..2147483647$ and \texttt{byte} is
$-128..127$ in Java, and Ada lets the programmer declare new ranges:
\begin{lstlisting}
type Population is range 0..1e10
type Row is range a..n
\end{lstlisting}
Composite types are built from simpler types by product, sum, mapping or
recursion. The lecture's examples are an Ada record, which is a product of
its fields, and a class hierarchy, which behaves as a disjoint union of the
classes that extend the base class:
\begin{lstlisting}
type Month is (jan, feb, ..., dec)
type Day is range 1..31
type Date is record
  m: Month;
  d: Day;
end record;
\end{lstlisting}

\subsection{Type constructors}

Table~\ref{tab:constructors} lists the constructors of the lecture. The names
$\Sigma$ and $\Pi$ come from Martin-L\"of's type theory, where they generalize
the cartesian product and the function space.

\begin{table}[h]
\centering
\small
\renewcommand{\arraystretch}{1.25}
\begin{tabular}{l l p{0.48\textwidth}}
  \toprule
  Constructor & Notation & Meaning \\
  \midrule
  Function   & $T_1 \to T_2$ & functions from $T_1$ to $T_2$ \\
  $\Sigma$-type & $\sum_{i=1}^{n} T_i = \{(v_1, \dots, v_n) \mid v_i : T_i\}$ & tuples of values of the given types; structures and records \\
  $\Pi$-type & $\prod_{v : T_1} T_2$ & functions whose result type may depend on the input value \\
  Equality   & $\mathit{exp}_1 : T = \mathit{exp}_2 : T$ & the type of proofs that two expressions reduce to the same one, as $2 + 2 = 2 * 2$ \\
  Inductive  & $0 : \mathbb{N}$,\ \ $s : \mathbb{N} \to \mathbb{N}$ & self-referential types built from constants and constructor functions \\
  \bottomrule
\end{tabular}
\caption{Type constructors.}
\label{tab:constructors}
\end{table}

\paragraph{Function types and currying.} The type of a function goes from the
types of its parameters to the type of its result. A function of several
parameters can always be seen as a function of one parameter that returns a
function: a function of two integers that returns an integer has type
$\mathit{int} \times \mathit{int} \to \mathit{int}$, or, curried,
$\mathit{int} \to (\mathit{int} \to \mathit{int})$. This is the currying property,
and the arrow associates to the right, so the parentheses are usually omitted:
$\mathit{int} \to \mathit{int} \to \mathit{int}$.

\paragraph{$\Sigma$ and $\Pi$.} A $\Sigma$-type is the type of tuples: if $T_1,
\dots, T_n$ are types, $\sum T_i$ contains the tuples $(v_1, \dots, v_n)$ with
$v_i : T_i$. A $\Pi$-type is the type of functions, where the type of the output
can depend on the input value: for each element $v$ of $T_1$ the function
returns an element of $T_2$, and $T_2$ may mention $v$. When $T_2$ does not
depend on $v$, $\prod_{v : T_1} T_2$ is just $T_1 \to T_2$. (In programming
languages the word ``sum type'' usually means a disjoint union, as
\texttt{Either} in Haskell; in type theory the disjoint union is written
$A + B$, and $\Sigma$ is a generalized product.)

\paragraph{Inductive types.} An inductive type is defined by constant elements
and constructor functions. The natural numbers of Peano have the constant
$0 : \mathbb{N}$ and the successor $s : \mathbb{N} \to \mathbb{N}$, and the
elements of the type are $0, s0, ss0, sss0, \dots$. Lists and trees are
inductive types too; Section~\ref{sec:recursive} returns to them.

\subsection{Static and dynamic type checking}

\paragraph{Static typing.} Every entity of the program has a type at compile
time, and the type checker verifies that the types are consistent before the
program runs. Its advantages: it catches errors early, it gives guarantees
about the programs, and the hardware can sometimes use the types directly
(an integer known at compile time needs no tag). Its cost: the programmer has
to write the types of the values, which type inference
(Section~\ref{sec:inference}) reduces.

\paragraph{Dynamic typing.} The same checks are made while the program runs.
For that, every value carries its type, as a pair (value, type). The lecture
uses Prolog to show it. A predicate
\texttt{list\_length([\_ | L], N)} computes the length of a list; called with a
list it gives $N = 4$, and called as \texttt{list\_length(9, N)} the tags say
that an \texttt{Int} is not a \texttt{List}, and the error is found only when
the call runs. The same idea in Scheme, with explicit tags:
\codelines{Scheme}{scheme/dynamic.scm}{2}{15}
The test calls \texttt{list-length} with a value tagged \texttt{int} and checks
that the error says \texttt{int is-not list}. Dynamic typing is more flexible
when programming, but execution costs more (the checks are repeated every
time an operation runs), the type information takes storage, and the lecture
adds that it promotes careless programming: an error in a branch that the
tests never execute stays hidden. Figure~\ref{fig:static-dynamic} contrasts
when each approach finds the error.

\begin{figure}[h]
\centering
\begin{tikzpicture}[box/.style={draw, minimum width=2.3cm, minimum height=7mm, font=\small}]
  \node[box] (s1) at (0,0) {source};
  \node[box] (s2) at (3.2,0) {type checker};
  \node[box] (s3) at (6.4,0) {execution};
  \draw[flecha] (s1) -- (s2);
  \draw[flecha] (s2) -- (s3);
  \node[font=\small, below=1mm of s2] {error found here (static)};
  \node[box] (d1) at (0,-1.8) {source};
  \node[box] (d3) at (6.4,-1.8) {execution};
  \draw[flecha] (d1) -- (d3);
  \node[font=\small, below=1mm of d3] {error found here, if the line runs (dynamic)};
  \node[font=\small, anchor=east] at (-1.4,0) {static};
  \node[font=\small, anchor=east] at (-1.4,-1.8) {dynamic};
\end{tikzpicture}
\caption{Where a type error is found.}
\label{fig:static-dynamic}
\end{figure}

\paragraph{Duck typing.} Types do not have to be structural: the type of an
expression can also be determined by its behavior. ``If it talks like a duck
and walks like a duck, then it is a duck.'' The lecture shows a Go interface
\texttt{Quacker} that any type with a method \texttt{Quack()} satisfies without
declaring it. In Python the check happens at run time:
\codelines{Python}{python/duck.py}{2}{15}
\texttt{make\_it\_quack} accepts a \texttt{Duck} and a \texttt{Robot}, which have
nothing in common except the method. With a \texttt{Dog} the error appears when
the call runs.

\subsection{Newer types}

Modern type systems add types that make common errors impossible.

\paragraph{Option types.} An option type (or Maybe type) $T?$ holds a value of
$T$ or nothing. It is used for function results that may have no meaningful
value. The lecture writes it in Scala, with \texttt{Some(42)} and
\texttt{None}; in Swift \texttt{Int?} is the same type, with \texttt{nil} for
nothing:
\codelines{Swift}{swift/optionals.swift}{2}{7}
To use the value the program has to open the option (\texttt{if let}), so the
case without a value cannot be forgotten.

\paragraph{Nullable types.} A \texttt{null} value makes operations unsafe: any
expression of a reference type may suddenly be \texttt{null}. A nullable type
system separates $T$ (always a valid reference) from $T?$ (a reference or
\texttt{null}), as in Figure~\ref{fig:nullable}, and the compiler checks that
a $T?$ is never used as a $T$ without a check.
\codelines{Swift}{swift/optionals.swift}{9}{18}
\texttt{greetUser} takes a \texttt{String}, not a \texttt{String?}, so the call
\texttt{greetUser(username: nil)} is rejected at compile time with
\texttt{'nil' is not compatible with expected argument type 'String'}.

\begin{figure}[h]
\centering
\begin{tikzpicture}[box/.style={draw, rounded corners=1pt, minimum height=7mm, font=\small}]
  \node[box] (opt) at (0,0) {\texttt{String?}};
  \node[box] (str) at (-1.6,-1.5) {\texttt{String}};
  \node[box] (nul) at (1.6,-1.5) {\texttt{Null}};
  \node[box, align=center] (obj) at (-1.6,-3.1) {object on the heap\\attributes, methods};
  \draw[flecha] (opt) -- node[left, font=\footnotesize] {may be} (str);
  \draw[flecha] (opt) -- node[right, font=\footnotesize] {may be} (nul);
  \draw[flecha] (str) -- node[left, font=\footnotesize] {valid reference} (obj);
\end{tikzpicture}
\caption{A nullable type is a valid reference or null.}
\label{fig:nullable}
\end{figure}

\paragraph{Owned types.} In Rust every value is attached to one variable, its
owner, and has exactly one owner at a time. Assigning a value or passing it as
a parameter moves it to a new owner, and the old variable can no longer be
used. A value can also be borrowed (\texttt{\&s}): passed temporarily to a
function and returned to its owner. In the lecture's example,
\texttt{print\_string(s)} moves the string \texttt{s} into the function, so a
second \texttt{print\_string(s)} is a compile error; \texttt{print\_string(\&s)}
only borrows it and can be repeated. Figure~\ref{fig:owned} shows what moves:
the variable holds a pointer, a length and a capacity, and the characters live
on the heap. Moving copies the three fields to the new owner and invalidates
the old one, so the heap memory always has a single owner that frees it.

\begin{figure}[h]
\centering
\begin{tikzpicture}[cell/.style={draw, minimum width=1.3cm, minimum height=6mm, font=\small\ttfamily}, hc/.style={draw, minimum width=7mm, minimum height=6mm, font=\small\ttfamily}]
  \node[font=\small] at (0,0.7) {\texttt{s1} (stack)};
  \node[cell] (p) at (0,0) {ptr};
  \node[cell] (l) at (0,-0.6) {len = 5};
  \node[cell] (c) at (0,-1.2) {cap = 5};
  \node[font=\small] at (4.5,0.7) {heap};
  \foreach \ch [count=\i from 0] in {h,e,l,l,o} {
    \node[hc] (h\i) at (3.1+\i*0.7,0) {\ch};
    \node[font=\scriptsize] at (3.1+\i*0.7,-0.5) {\i};
  }
  \draw[flecha] (p.east) -- (h0.west);
  \node[font=\small, align=center] at (0,-2.2) {after \texttt{takes\_ownership(s1)}\\the fields belong to the function};
\end{tikzpicture}
\caption{A Rust \texttt{String}: the owner holds pointer, length and capacity; the text is on the heap.}
\label{fig:owned}
\end{figure}

\subsection{Well-typed programs cannot go wrong}

Robin Milner summarized the goal of static type systems as ``well-typed
programs cannot go wrong''. Going wrong has a precise meaning: reaching a
state where no rule applies, such as applying the number 3 as a function. The
claim is proved as two theorems, progress and preservation, in
Section~\ref{sec:safety}. The sections in between build what the proof needs:
typing rules and the language LETREC (Section~\ref{sec:letrec}), the lambda
calculus and its typed version, type inference, polymorphism, subtyping and
recursive types.

\subsection{Exercises}

\begin{exercise}
Classify Python, Java and Swift by the four questions of this section
(static or dynamic, strong or weak, explicit or inferred, nominal or
structural).
\end{exercise}

\paragraph{Solution.} Python: dynamic, strong, no declarations (types are
checked at run time), and duck typing, which is structural by behavior. Java:
static, strong, explicit (with local inference through \texttt{var}), nominal:
a class is a subtype of an interface only if it declares \texttt{implements}.
Swift: static, strong, inferred where possible, nominal for classes and
protocols (a type conforms to a protocol by declaring it).

\begin{exercise}
Give the type of a curried function \texttt{add} of two integers, and the type
of \texttt{add 3}.
\end{exercise}

\paragraph{Solution.} $\mathit{add} : \mathit{int} \to \mathit{int} \to \mathit{int}$,
which is $\mathit{int} \to (\mathit{int} \to \mathit{int})$. Applying it to 3 gives a
function $\mathit{int} \to \mathit{int}$ that adds 3.

\begin{exercise}
Write the natural number 3 and the list $[1, 2]$ as elements of inductive
types.
\end{exercise}

\paragraph{Solution.} $3 = s(s(s(0)))$. With the constructors
$\mathit{nil} : \mathit{List}$ and $\mathit{cons} : \mathbb{N} \to \mathit{List} \to \mathit{List}$,
$[1, 2] = \mathit{cons}(s(0), \mathit{cons}(s(s(0)), \mathit{nil}))$.

\begin{exercise}
Does ``well-typed programs cannot go wrong'' mean that well-typed programs
always terminate?
\end{exercise}

\paragraph{Solution.} No. A well-typed program never gets stuck, but it may run
forever: in LETREC, \texttt{letrec int f (x : int) = (f x) in (f 0)} is well
typed and loops.
